\section{Cycle Equivalence}

Subsection~3.1 defines \texttt{Cycle} to mean \texttt{RecCycle} --- a naming
choice, not a claim. This section proves the claim that actually needs proving:
that \texttt{ModCycle}, defined independently by a single modulo operation
rather than by recursive unrolling, computes exactly the same values as
\texttt{RecCycle} at every position. This equivalence is what the rest of the
article leans on whenever it moves between the two representations, and it is
not immediate from the two definitions side by side. We prove it by induction
on the position $i$.

\begin{itemize}
  \item Base case ($i<n$): both definitions consult the base list directly at
    the same position.
  \item Inductive step ($i\geq n$): the recursive definition reduces to $i-n$,
    wrapping to the same modulo position.
\end{itemize}

\begin{equation*}
\begin{aligned}
\forall\, L\in\mathbb{L},\quad \forall\, v&\in\mathbb{N}_0,
  \quad \forall\, i\in\mathbb{N}_0 \\
L &:= [v_0,v_1,\dots,v_{n-1}]\in\mathbb{N}_0^n \\
n &= |L| \\
\operatorname{ModCycle}_i &= L[i\operatorname{mod}n],\quad |L|>0 \\
\operatorname{RecCycle}_i &=
\begin{cases}
L_i & \text{if } i<n \\
\operatorname{RecCycle}_{i-n} & \text{if } i\geq n
\end{cases},\quad |L|>0
\end{aligned}
\end{equation*}

\subsection{Base Case (\texorpdfstring{$i<n$}{i < n})}

When the position is within the first cycle, both definitions return the list
element directly.

\begin{equation*}
\begin{aligned}
i<n \implies i\operatorname{mod}n &= i
  && \text{[Trivial Mod For Small Dividend]} \\
\operatorname{ModCycle}_i &= L_{(i\operatorname{mod}n)}
  && \text{[ModCycle Definition]} \\
&=L_i && \text{[Since }i<n\text{, }i\operatorname{mod}n=i\text{]} \\
\operatorname{RecCycle}_i &= L_i
  && \text{[Since }i<n\text{, by RecCycle Definition]} \\
&\therefore \\
i<n \implies \operatorname{ModCycle}_i &= \operatorname{RecCycle}_i
  \quad \blacksquare && \text{[Q.E.D.]}
\end{aligned}
\end{equation*}

The lemma
\href{https://github.com/thiagomata/prime-numbers/blob/cycle-article-v1.0.0/articles/chapter2/modulo.md#61-trivial-case}{Trivial Mod for Small Dividend}
was proved and verified in the article
\href{https://github.com/thiagomata/prime-numbers/blob/cycle-article-v1.0.0/articles/chapter2/modulo.md}{\emph{Division and Modulo from Recursive Normalization}}~\cite{mata2026modulo}.

{\raggedright
This property is verified in
\href{https://github.com/thiagomata/prime-numbers/blob/cycle-article-v1.0.0/src/main/scala/v1/chapter4/cycle/recursive/properties/RecursiveCycleMatchesModCycle.scala}{\texttt{RecursiveCycleMatchesModCycle::}\allowbreak\texttt{assertCycleAndRecursiveCycleMathForSmallValues}}.
The full Scala verification code is in Appendix~A.1.
\par}

\subsection{Inductive Step (\texorpdfstring{$i\geq n$}{i >= n})}

For positions beyond the first cycle, both definitions reduce the position by
the cycle period and rely on the inductive hypothesis.

\begin{equation*}
\begin{aligned}
\operatorname{ModCycle}_{(i-n)} &= \operatorname{RecCycle}_{(i-n)}
  && \text{[By Induction Hypothesis]} \\
i\geq n \implies i\operatorname{mod}n
  &= (i-n)\operatorname{mod}n
  && \text{[Quotient Invariance Under Linear Shift]} \\
\operatorname{ModCycle}_i &= L_{(i\operatorname{mod}n)}
  && \text{[ModCycle Definition]} \\
&=L_{((i-n)\operatorname{mod}n)}
  && \substack{\text{[Since }i\geq n\text{, }i\operatorname{mod}n
       =(i-n)\operatorname{mod}n\text{]}} \\
&=\operatorname{ModCycle}_{(i-n)} && \text{[By Definition]} \\
&=\operatorname{RecCycle}_{(i-n)} && \text{[By Substitution]} \\
\operatorname{RecCycle}_i &= \operatorname{RecCycle}_{(i-n)}
  && \text{[By RecCycle Definition]} \\
&=\operatorname{ModCycle}_i && \text{[By Substitution]} \\
&\therefore \\
i\geq n \implies \operatorname{ModCycle}_i &= \operatorname{RecCycle}_i
  \quad \blacksquare && \text{[Q.E.D.]}
\end{aligned}
\end{equation*}

The lemma
\href{https://github.com/thiagomata/prime-numbers/blob/cycle-article-v1.0.0/articles/chapter2/modulo.md#65-quotient-invariance-under-linear-shift}{Quotient Invariance Under Linear Shift}
was proved and verified in
\href{https://github.com/thiagomata/prime-numbers/blob/cycle-article-v1.0.0/articles/chapter2/modulo.md}{\emph{Division and Modulo from Recursive Normalization}}~\cite{mata2026modulo}.

{\raggedright
This property is verified in
\href{https://github.com/thiagomata/prime-numbers/blob/cycle-article-v1.0.0/src/main/scala/v1/chapter4/cycle/recursive/properties/RecursiveCycleMatchesModCycle.scala}{\texttt{RecursiveCycleMatchesModCycle::}\allowbreak\texttt{assertCycleAndRecursiveCycleMathForAnyValues}}.
The full Scala verification code is in Appendix~A.2.
\par}

With both pieces now established --- \texttt{RecCycle} and \texttt{ModCycle}
equal by this section, \texttt{MemCycle} and \texttt{ModCycle} equal by
delegation (Subsection~3.3) --- the three representations agree at every
position:

\begin{equation*}
\begin{aligned}
\operatorname{RecCycle}(L)_i
  &=\operatorname{ModCycle}(L)_i && \text{[Subsections 4.1--4.2]} \\
  &=\operatorname{MemCycle}(L)_i && \text{[Subsection 3.3]} \\
\therefore\quad \operatorname{RecCycle}(L)_i
  &=\operatorname{ModCycle}(L)_i=\operatorname{MemCycle}(L)_i
  && \text{[Three-Way Equality]}
\end{aligned}
\end{equation*}
